\section{Work Execution Timeline}

The work will be executed over six months, organized into four main stages aligned with the work description. The timeline accounts for literature review, implementation, experimentation, analysis, and thesis writing.

\begin{itemize}[leftmargin=0pt]
    \item \textbf{Month 1:} \textbf{Stage 1 - Baseline Comparison and Validation}
    \begin{itemize}
        \item Literature review and analysis of baseline methods (DeepTTE~\cite{deepTTE2018}, STAD~\cite{abbar2020stad}, DCRNN~\cite{dcrnn2018}, ST-GCN~\cite{yu2018stgcn}, DuETA~\cite{dueta2023})
        \item Dataset preparation and preprocessing: chronological partitioning, feature normalization, temporal window construction
        \item Implementation of baseline models and evaluation framework
        \item Initial baseline experiments and performance benchmarking
    \end{itemize}
    
    \item \textbf{Month 2:} \textbf{Stage 1 Completion and Stage 2 Initiation}
    \begin{itemize}
        \item Completion of baseline comparison experiments
        \item Analysis of baseline results and validation of dataset suitability
        \item Begin Stage 2: Implementation of route-aware and route-agnostic DSTRA-GNN variants
        \item Design and implementation of ablation study framework
    \end{itemize}
    
    \item \textbf{Month 3:} \textbf{Stage 2 - Route Information Contribution Analysis}
    \begin{itemize}
        \item Systematic ablation studies comparing route-aware vs. route-agnostic variants
        \item Experiments across different trip categories and traffic conditions
        \item Analysis of route information contribution across scenarios
        \item Statistical analysis and significance testing of results
    \end{itemize}
    
    \item \textbf{Month 4:} \textbf{Stage 3 - Architectural Enhancements}
    \begin{itemize}
        \item Design and implementation of enhanced route encoding mechanisms
        \item Development of improved temporal aggregation strategies
        \item Refinement of Mixture-of-Experts routing mechanism
        \item Integration of additional contextual features (if applicable)
        \item Training and evaluation of enhanced architecture variants
    \end{itemize}
    
    \item \textbf{Month 5:} \textbf{Stage 4 - Comprehensive Evaluation and Analysis}
    \begin{itemize}
        \item Final model training and hyperparameter tuning
        \item Comprehensive evaluation: performance comparison, error analysis, computational efficiency assessment
        \item Visualization of learned representations and interpretability analysis
        \item Statistical analysis and synthesis of all experimental results
        \item Identification of failure modes and limitations
    \end{itemize}
    
    \item \textbf{Month 6:} \textbf{Thesis Writing and Finalization}
    \begin{itemize}
        \item Writing thesis chapters: introduction, methodology, experiments, results, discussion, conclusion
        \item Preparation of figures, tables, and visualizations
        \item Final review and refinement of experimental results
        \item Thesis editing, proofreading, and formatting
        \item Final submission preparation
    \end{itemize}
\end{itemize}

\textbf{Note:} The timeline allows for iterative refinement, where insights from later stages may inform adjustments to earlier components. Regular progress reviews with the supervisor will ensure alignment with research objectives and timely completion.

